% Intended LaTeX compiler: pdflatex
\documentclass[a4paper,12pt]{article}
\usepackage[utf8]{inputenc}
\usepackage[T1]{fontenc}
\usepackage{graphicx}
\usepackage{longtable}
\usepackage{wrapfig}
\usepackage{rotating}
\usepackage[normalem]{ulem}
\usepackage{amsmath}
\usepackage{amssymb}
\usepackage{capt-of}
\usepackage{hyperref}
\usepackage{\string~/.emacs.d/common}
\author{mahmood}
\date{\today}
\title{algorithm correctness}
\hypersetup{
 pdfauthor={mahmood},
 pdftitle={algorithm correctness},
 pdfkeywords={},
 pdfsubject={},
 pdfcreator={Emacs 30.0.50 (Org mode 9.6)}, 
 pdflang={English}}
\begin{document}

\maketitle
\begin{definition}
an algorithm is totally correct if it receives valid input, terminates, and always returns the correct output\\[0pt]
\end{definition}
the only way to prove the \href{20221104220603-algorithm_correctness.org}{correctness of an algorithm} over all possible inputs is by reasoning formally or mathematically about it.\\[0pt]
generally, to prove an algorithm is correct we prove two things\\[0pt]
\begin{enumerate}
\item the algorithm terminates\\[0pt]
\item the algorithm returns the correct value for every input\\[0pt]
\end{enumerate}
we propose a statement that describes what the algorithm takes as input and what its supposed to output and then prove it\\[0pt]
\section{proving algorithm correctness by induction}
\label{sec:org3c5e520}
here we'll be using \href{20220707193301-mathematical_induction.org}{mathematical induction}\\[0pt]
\begin{my_example}
\includesvg{/Users/mahmooz/brain/out/cueVWTf}\\[0pt]

\begin{lemma}
for each input that is an array of numbers A, the algorithm \(\textsc{Sum-Array}\) computes the sum of the numbers\\[0pt]
\begin{proof}
\uline{proof the algorithm terminates}: the algorithm consists of a for loop that is bound by n, so it follows trivially that it must terminate\\[0pt]
\uline{proof of correctness} by \href{20220707193301-mathematical_induction.org}{perfect induction} on the length of the input array\\[0pt]
we propose the following \href{20221104221055-loop_invariant.org}{loop invariant} for the loop\\[0pt]
\[
sum = \sum_{j=1}^{i} A[j]
\]\\[0pt]
in this case, proving the proposed loop invariant would prove the algorithm correct\\[0pt]
\begin{enumerate}
\item for \(i=1\), which is the base case, the algorithm returns A[1] which would be the only element in the array, so its obvious that its the correct output\\[0pt]
\item the \href{20230108230724-induction_hypothesis.org}{induction hypothesis}, we assume after i iterations in the loop that \(sum = \sum_{j=1}^{i} A[j]\)\\[0pt]
\item the \href{20230108230748-induction_step.org}{induction step}, during the \(i+1\)'th iteration of the loop, we add \(A[i+1]\) to the variable sum, and so it follows from this and from the induction hypothesis that:\\[0pt]
\end{enumerate}
\[
sum = \sum_{j=1}^{i} A[j] + A[i+1] = \sum_{j=1}^{i+1} A[j]
\]\\[0pt]
the loop runs n times, by substituting n in the loop invariant, we'd get\\[0pt]
\[
sum = \sum_{j=1}^{n} A[j]
\]\\[0pt]
which means that after the loop is done, the variable sum would contain the sum of the elements in the array, which is what the function returns, which is the correct output\\[0pt]
QED\\[0pt]
\end{proof}
\end{lemma}
\end{my_example}
\section{correctness of recursive algorithms}
\label{sec:orgdce3b3c}
the pattern for proofs of recursive algorithms is as follows:\\[0pt]
\begin{enumerate}
\item \uline{base step of the recursion}, we have to manually show that for each of the edge cases the output is correct\\[0pt]
\item \uline{induction hypothesis}, we assume that in each recurring call the output is correct\\[0pt]
\item \uline{proof}, we show that the output is correct for the main function call\\[0pt]
\end{enumerate}
\begin{my_example}
consider the following function\\[0pt]
\includesvg{/Users/mahmooz/brain/out/e8I9jOY}\\[0pt]

\uline{proof that the algorithm terminates}:\\[0pt]
assume that \(\textsc{Recursive-Sum}\) runs on an array \texttt{A}, on each recurring call, the size of the input array \texttt{A} is halved, its obvious that at the end of every sequence of calls the size of the input array would arrive to 1 which is the edge case that would cause the function to terminate\\[0pt]
\uline{proof of correctness}:\\[0pt]
we prove this by induction on the number of elements in the array\\[0pt]
\begin{enumerate}
\item \uline{base case}: for \(|A|=1\) which is the edge case, the algorithm returns \(A[1]\), the only element in \(A\), its obvious that the output is correct\\[0pt]
\item \uline{induction hypothesis}: we use \href{20220707193301-mathematical_induction.org}{perfect induction} on the size of the input, we assume that the algorithm correctly computes the sum of the elements of any array whose size is smaller than \texttt{n}\\[0pt]
\item \uline{induction step}: let \texttt{A} be an arbitrary array of size \texttt{n}, the call to \(\textsc{Recursive-Sum}(A)\) executes the calls \(\textsc{Recursive-Sum}(A_l)\) and \(\textsc{Recursive-Sum}(A_r)\), the size of the arrays \(A_l\) and \(A_r\) is smaller than \(n\), so according to the previous step we know those two recurring calls do return the correct sum of the corresponding subarrays, and since we are returning the sum of those sums we know for sure we are returning the correct sum\\[0pt]
\end{enumerate}
\end{my_example}
\end{document}